\section{Projective resolutions and Ext groups of standard modules}
\label{sec:projective-resolutions-ext-groups-standard-modules}

In this section, we study projective resolutions of standard
$\Lambda_P$-modules and compute Ext groups between standard and simple modules.

\subsection{Projective resolutions of standard modules}
\label{sec:projective-resolutions-standard-modules}

For every $T\in\Int(P)$, we construct an explicit projective
resolution of $\DD(T)$ from the combinatorics of relative upsets of
$T$ (\Cref{thm:standard-projective-resolution}). 

\subsubsection{Combinatorial preliminaries}
\label{subsec:projective-resolution-combinatorics}

Fix an interval $T\in\Int(P)$.  
Let $\cU^{\circ}(T):=\cU(T)\setminus\{T\}$ be the set of intervals contained in $T$ that are proper relative upsets in $T$. 
Choose an ordering $U_1,\ldots,U_m$ of the inclusion-maximal elements
of $\cU^\circ(T)$ and write
\begin{equation}
\cU^\circ_{\max}(T)=\{U_1,\ldots,U_m\}
\qand
m=|\cU^\circ_{\max}(T)|.
\label{eq:maximal-proper-relative-upsets}
\end{equation}
For $J\subseteq[m]$, put
\[
U_J=\bigcap_{j\in J}U_j
\]
with the convention $U_\varnothing=T$. 
If $U_J=\varnothing$, it contributes no
connected components.

Let $J'\subseteq J\subseteq[m]$. Since $U_J\subseteq U_{J'}$, every
$C\in\piZero(U_J)$ is contained in a unique connected component of
$U_{J'}$. We denote this component by
\[
C^{J'}\in\piZero(U_{J'}).
\]

\begin{lemma}
\label{lem:projective-resolution-face-components}
Let $J'\subseteq J\subseteq[m]$ and
$C\in\piZero(U_J)$. Then $C\in \cU(C^{J'})$.
Moreover,
\begin{equation}
\label{eq:face-component-functoriality}
\bigl(C^{J'}\bigr)^{J''}
=
C^{J''}
\qquad
(J''\subseteq J'\subseteq J).
\end{equation}
\end{lemma}

\begin{proof}
Since $U_J$ is a relative upset of $U_{J'}$, each connected component
$C$ of $U_J$ is a relative upset of the component $C^{J'}$ of
$U_{J'}$ containing it.  Hence $C\in\cU(C^{J'})$.

The second assertion follows from the uniqueness of the connected
component containing $C$.
\end{proof}

For $D\in\cU^\circ(T)$, put
\begin{equation}
I_T(D)
=
\bigl\{
i\in[m]
\bigm|
D\subseteq U_i
\bigr\}.
\label{eq:ITD-definition}
\end{equation}
Every such $D$ is contained in some
$U_i\in\cU^\circ_{\max}(T)$, and hence $I_T(D)$ is nonempty.

The following lemma records the condition under which $D$ is contained
in $U_J$.

\begin{lemma}
\label{lem:unique-lift-admissible-support}
Let $D\in\cU^\circ(T)$. For $J\subseteq[m]$,
the following conditions are equivalent:
\begin{enumerate}[\rm (1)]
    \item $J\subseteq I_T(D)$;
    \item $D\subseteq U_J$;
    \item there exists a unique $C\in\piZero(U_J)$ such that
    $D\subseteq C$.
\end{enumerate}
\end{lemma}

\begin{proof}
The equivalence of \textup{(1)} and \textup{(2)} follows directly from
the definitions. If \textup{(2)} holds, then the connected subset $D$
is contained in a unique connected component of $U_J$, proving
\textup{(3)}. Conversely, \textup{(3)} gives
$D\subseteq C\subseteq U_J$, and hence \textup{(2)}.
\end{proof}

Consequently, the subsets $J\subseteq[m]$ satisfying
$D\subseteq U_J$ form the full simplex with vertex set $I_T(D)$:
\begin{equation}
\bigl\{
J\subseteq[m]
\bigm|
D\subseteq U_J
\bigr\}
=
\Simp(I_T(D)).
\label{eq:ITD-simplex}
\end{equation}


For $D\in\cU^\circ(T)$ and
$J\subseteq I_T(D)$, denote the unique component in
\Cref{lem:unique-lift-admissible-support} by
\begin{equation}
\label{eq:compJD-definition}
\operatorname{comp}_J(D)\in\piZero(U_J).
\end{equation}
By construction,
$D\subseteq\operatorname{comp}_J(D)$ and
$\operatorname{comp}_\varnothing(D)=T$.
Moreover, \Cref{lem:projective-resolution-face-components}, applied
with $J'=\varnothing$, gives
$\operatorname{comp}_J(D)\in\cU(T)$.
The same uniqueness also gives
\begin{equation}
\label{eq:lifted-component-compatibility}
\bigl(\operatorname{comp}_J(D)\bigr)^{J'}
=
\operatorname{comp}_{J'}(D)
\qquad
\bigl(J'\subseteq J\subseteq I_T(D)\bigr).
\end{equation}


\begin{lemma}
\label{lem:projective-resolution-admissibility-transfer}
Let $R\in\Int(P)$ and $C\in\cU(T)$. Then
\begin{equation}
\label{eq:admissibility-transfer}
\Adm(R,C)
=
\bigl\{
D\in\Adm(R,T)
\bigm|
D\subseteq C
\bigr\}.
\end{equation}
\end{lemma}

\begin{proof}
Let $D\in\Adm(R,C)$. Then
$D\in\cD(R)$ and $D\in\cU(C)$.
Since $C\in\cU(T)$, we have
$D\in\cU(T)$ and $D\subseteq C$.
Thus $D$ belongs to the right-hand side of
\eqref{eq:admissibility-transfer}.

Conversely, let $D$ belong to the right-hand side of
\eqref{eq:admissibility-transfer}. Then
$D\in\cD(R)$, $D\in\cU(T)$, and $D\subseteq C$.
Since $C\in \cU(T)$, we have $D\in\cU(C)$. Thus $D\in\Adm(R,C)$.
\end{proof}

For $R\in\Int(P)$, write
\begin{equation}
\label{eq:proper-admissible-supports}
\Adm^\circ(R,T)
=
\Adm(R,T)\setminus\{T\}.
\end{equation}
For any $D\in\Adm^\circ(R,T)$, we have $D \in \cU^\circ(T)$, and hence $I_T(D)$ is nonempty.
Moreover, for each $J\subseteq I_T(D)$, one has
$D\subseteq\operatorname{comp}_J(D)$ and
$\operatorname{comp}_J(D)\in\cU(T)$.
Therefore, \eqref{eq:admissibility-transfer} yields
$D\in\Adm\bigl(R,\operatorname{comp}_J(D)\bigr)$.

Consequently, for every $R\in\Int(P)$ and every nonempty
$J\subseteq[m]$, projection onto the second factor induces a bijection
\begin{equation}
\label{eq:admissible-support-component-bijection}
\bigsqcup_{C\in\piZero(U_J)}
\Adm(R,C)
\xrightarrow{\sim}
\bigl\{
D\in\Adm^\circ(R,T)
\bigm|
J\subseteq I_T(D)
\bigr\},
\qquad
(C,D)\longmapsto D.
\end{equation}
Its inverse sends $D$ to
$\bigl(\operatorname{comp}_J(D),D\bigr)$.

\subsubsection{Construction of projective resolutions}
\label{subsec:projective-resolution-construction}
We now construct a complex $\QQ_\bullet(T)$ of projective
$\Lambda_P$-modules for $T\in\Int(P)$, indexed by the integers
$q\geq0$.  
Let $\cU^\circ_{\max}(T)=\{U_1,\ldots,U_m\}$ be as in
\eqref{eq:maximal-proper-relative-upsets}.
For $q\geq0$, set
\begin{equation}
\label{eq:explicit-projective-terms}
\QQ_q(T)
=
\bigoplus_{\substack{J\subseteq[m]\\ |J|=q}}
\ \bigoplus_{C\in\piZero(U_J)}
\PP(C).
\end{equation}
Since $U_\varnothing=T$ and $T$ is connected, this gives
$\QQ_0(T)=\PP(T)$.

Let $q\geq1$, let
$J=\{j_0<\cdots<j_{q-1}\}\subseteq[m]$,
and let $C\in\piZero(U_J)$. For $0\leq r\leq q-1$, put
\[
J_r=J\setminus\{j_r\}.
\]
By \Cref{lem:projective-resolution-face-components},
$C\in\cU(C^{J_r})$. Define
\[
\partial_{J,C}^{r}
\colon
\PP(C)\longrightarrow\QQ_{q-1}(T)
\]
to be the composite of
\[
(j_{C,C^{J_r}})_*
\colon
\PP(C)\longrightarrow\PP(C^{J_r})
\]
with the canonical inclusion of the summand of $\QQ_{q-1}(T)$ indexed
by $(J_r,C^{J_r})$. Set
\begin{equation}
\label{eq:explicit-projective-component-differential}
d_{J,C}
=
\sum_{r=0}^{q-1}(-1)^r\partial_{J,C}^{r}
\colon
\PP(C)\longrightarrow\QQ_{q-1}(T).
\end{equation}
As $(J,C)$ ranges over the summands in
\eqref{eq:explicit-projective-terms}, the maps $d_{J,C}$ define a
homomorphism
\begin{equation}
\label{eq:explicit-projective-differential}
d_q
\colon
\QQ_q(T)\longrightarrow\QQ_{q-1}(T)
\qquad
(q\geq1).
\end{equation}
In particular,
$\QQ_1(T)=\bigoplus_{i=1}^{m}\PP(U_i)$, and $d_1$ is the sum of the
maps
\[
(j_{U_i,T})_*\colon\PP(U_i)\longrightarrow\PP(T).
\]

Every summand of $d_{q-1}d_q$ is obtained by deleting two indices from
$J$.  By \eqref{eq:face-component-functoriality} and the inclusion-map
composition law in \Cref{lem:interval-directional-composition}, the two
possible orders of deletion contribute the same projective map with
opposite signs.  Hence
\begin{equation}
\label{eq:explicit-projective-complex-condition}
d_{q-1}d_q=0
\qquad
(q\geq2).
\end{equation}
Thus $\QQ_\bullet(T)$ is a complex of projective
$\Lambda_P$-modules.

The following result shows that adjoining the canonical epimorphism
$\pi_T$ to $\QQ_\bullet(T)$ gives a projective resolution of $\DD(T)$.

\begin{theorem}
\label{thm:standard-projective-resolution}
For every interval $T\in\Int(P)$, the sequence
\begin{equation}
\label{eq:standard-projective-resolution-sequence}
\cdots
\longrightarrow
\QQ_2(T)
\xrightarrow{d_2}
\QQ_1(T)
\xrightarrow{d_1}
\QQ_0(T)
\xrightarrow{\pi_T}
\DD(T)
\longrightarrow0
\end{equation}
is exact. Consequently, it is a projective resolution of
$\DD(T)$.
\end{theorem}


In the rest of this subsubsection, we prove
\Cref{thm:standard-projective-resolution}.
Since exactness of
\eqref{eq:standard-projective-resolution-sequence}
is detected componentwise, it suffices to prove that
\begin{equation}
\label{eq:eR-standard-projective-resolution-sequence}
\cdots
\longrightarrow
e_R\QQ_2(T)
\xrightarrow{e_Rd_2}
e_R\QQ_1(T)
\xrightarrow{e_Rd_1}
e_R\QQ_0(T)
\xrightarrow{e_R\pi_T}
e_R\DD(T)
\longrightarrow0
\end{equation}
is exact for every $R\in\Int(P)$.

Fix $R\in\Int(P)$. For $D\in\Adm^\circ(R,T)$ and
$J\subseteq I_T(D)$,
\eqref{eq:admissibility-transfer} gives
$D\in\Adm\bigl(R,\operatorname{comp}_J(D)\bigr)$. Set
\begin{equation}
\label{eq:xiJD-definition}
\xi_{J,D}
:=
\rho_{R,\operatorname{comp}_J(D)}^{D}
\in
e_R\PP\bigl(\operatorname{comp}_J(D)\bigr).
\end{equation}
In particular, 
$\xi_{\varnothing,D}=\rho_{R,T}^{D}
\in e_R\QQ_0(T)$.


With this notation, for every $q\geq1$, we have
\begin{align}
e_R\QQ_q(T)
&\overset{\eqref{eq:explicit-projective-terms}}{=}
\bigoplus_{\substack{J\subseteq[m]\\|J|=q}}
\ \bigoplus_{C\in\piZero(U_J)}
e_R\PP(C)
\notag\\
&\overset{\eqref{eq:projective-injective-component-bases}}{=}
\bigoplus_{\substack{J\subseteq[m]\\|J|=q}}
\ \bigoplus_{C\in\piZero(U_J)}
\ \bigoplus_{D\in\Adm(R,C)}
\kk\rho_{R,C}^{D}
\notag\\
&\overset{\eqref{eq:admissible-support-component-bijection}}{=}
\bigoplus_{D\in\Adm^\circ(R,T)}
\ \bigoplus_{\substack{J\subseteq I_T(D)\\|J|=q}}
\kk\xi_{J,D}.
\label{eq:eR-projective-terms-support-decomposition}
\end{align}

The decomposition
\eqref{eq:eR-projective-terms-support-decomposition}
groups the basis vectors according to
$D\in\Adm^\circ(R,T)$.  
For such a $D$, let
\begin{equation}
\label{eq:VRTD-definition}
\VV_q^{R,T}(D)
=
\bigoplus_{\substack{J\subseteq I_T(D)\\|J|=q}}
\kk\xi_{J,D}
\qquad
(q\geq0).
\end{equation}
In particular,
$\VV_0^{R,T}(D)=\kk\xi_{\varnothing,D} = \kk\rho_{R,T}^D$.

We claim that the differential $e_Rd_\bullet$ preserves these
summands and restricts to
$\VV_\bullet^{R,T}(D)$, making it a subcomplex of
$e_R\QQ_\bullet(T)$. 
Indeed, let $q\geq1$ and let
$J=\{j_0<\cdots<j_{q-1}\}\subseteq I_T(D)$.  Put
$J_r=J\setminus\{j_r\}$.  By
\eqref{eq:lifted-component-compatibility},
\eqref{eq:explicit-projective-component-differential}, and
\eqref{eq:admissible-basis-composition},
\begin{equation}
\label{eq:admissible-support-differential}
(e_Rd_q)(\xi_{J,D})
=
\sum_{r=0}^{q-1}(-1)^r\xi_{J_r,D}.
\end{equation}
Thus the differential preserves
$\VV_\bullet^{R,T}(D)$, proving the claim.


\begin{lemma}
\label{lem:admissible-support-simplex-complex}
For every $D\in\Adm^\circ(R,T)$, the complex
$\VV_\bullet^{R,T}(D)$ is a direct summand of
$e_R\QQ_\bullet(T)$.
Moreover,
\begin{equation}
\VV_\bullet^{R,T}(D)
\cong
\widetilde C_\bullet
\bigl(\Simp(I_T(D));\kk\bigr)[1]
\label{eq:VRTD-simplex-isomorphism}
\end{equation}
as chain complexes. 
In particular, $\VV_\bullet^{R,T}(D)$ is exact.
\end{lemma}

\begin{proof}
In each degree, $\VV_{q}^{R,T}(D)$ is spanned by the basis vectors
whose second index is $D$.  By
\eqref{eq:admissible-support-differential}, the differential of such
a basis vector is a linear combination of basis vectors with the same
second index.  Moreover, a basis vector indexed by $D'\neq D$ has no
component indexed by $D$ after applying the differential.  Hence the
degreewise projection onto the basis vectors indexed by $D$ is a
chain map and restricts to the identity on
$\VV_{\bullet}^{R,T}(D)$.  Thus $\VV_{\bullet}^{R,T}(D)$ is a direct summand
subcomplex of $e_R\QQ_\bullet(T)$.

For $q\geq1$, the basis vectors of $\VV_q^{R,T}(D)$ are indexed by the
$q$-element subsets of $I_T(D)$, while $\VV_0^{R,T}(D)$ is spanned by
$\rho_{R,T}^{D}$.  Hence there are natural identifications
\[
\VV_q^{R,T}(D)
\cong
\widetilde C_{q-1}\bigl(\Simp(I_T(D));\kk\bigr),
\qquad q\geq0,
\]
where $\rho_{R,T}^{D}$ corresponds to the augmentation term in degree
$-1$.  
Under these identifications,
\eqref{eq:admissible-support-differential} is the usual alternating
simplicial boundary. After multiplying the identification in degree
$q$ by $(-1)^q$, this gives the chain isomorphism
\eqref{eq:VRTD-simplex-isomorphism}.
\end{proof}

We now complete the proof of
\Cref{thm:standard-projective-resolution}.


\begin{proof}[Proof of \Cref{thm:standard-projective-resolution}]
Fix $R\in\Int(P)$. 
By \eqref{eq:eR-projective-terms-support-decomposition} and the
admissible-component basis of $e_R\PP(T)$, the complex
$e_R\QQ_\bullet(T)$ decomposes as
\begin{equation}
e_R\QQ_\bullet(T)
=
\left(
\bigoplus_{D\in\Adm^\circ(R,T)}
\VV_\bullet^{R,T}(D)
\right)
\oplus
\begin{cases}
\kk q_{R,T}[0]
& \text{if } R\in\cE_T,\\
0
& \text{otherwise}.
\end{cases}
\label{eq:eR-projective-complex-support-decomposition}
\end{equation}


By \eqref{eq:standard-components-explicit}, if $R\in\cE_T$, then
\begin{equation}
e_R\pi_T\colon
\kk q_{R,T}
\xrightarrow{\sim}
e_R\DD(T),
\qquad
q_{R,T}\longmapsto[q_{R,T}]
\label{eq:eR-standard-augmentation-complement}
\end{equation}
is an isomorphism.
We have $e_R\DD(T)=0$ otherwise.

Hence the augmented complex
\eqref{eq:eR-standard-projective-resolution-sequence} decomposes into
the complexes $\VV_\bullet^{R,T}(D)$, with
$D\in\Adm^\circ(R,T)$, and, when $R\in\cE_T$, the two-term complex
in \eqref{eq:eR-standard-augmentation-complement}.
By \Cref{lem:admissible-support-simplex-complex}, all these complexes
are exact. Therefore
\eqref{eq:eR-standard-projective-resolution-sequence} is exact.
\end{proof}


The projective resolution in
\Cref{thm:standard-projective-resolution} need not be minimal in general.





